\documentclass[12pt,a4paper]{article}
\usepackage[UTF8,heading=true]{ctex}
\usepackage{geometry}
\usepackage{amsmath,amssymb,amsthm}
\usepackage{graphicx}
\usepackage{booktabs}
\usepackage{longtable}
\usepackage{array}
\usepackage{enumitem}
\usepackage{hyperref}
\usepackage{fancyhdr}
\usepackage{xcolor}
\usepackage{listings}
\usepackage{setspace}  % 添加行距控制包

% 配置 ctex 章节格式为 "一、"
\ctexset{
  section = {
    name = {,、},
    number = \chinese{section},
    format = {\Large\bfseries\raggedright}
  }
}

\geometry{left=2.5cm,right=2.5cm,top=2.5cm,bottom=2.5cm}
\setlength{\headheight}{15pt}
\setstretch{1.5}  % 设置1.5倍行距（CNIPA推荐）

\pagestyle{fancy}
\fancyhf{}
\fancyhead[C]{发明专利技术交底书}
\fancyfoot[C]{\thepage}

\title{\textbf{发明专利技术交底书} \\[0.5em]
\Large 基于注意力语义感知的静默KV缓存淘汰与原位压缩的方法、装置、芯片及介质}
\author{中国移动通信有限公司研究院}
\date{\today}

\begin{document}

\begin{center}
    {\Large\heiti 中国移动专利申请} \\[1em]
    {\Huge\heiti 技术交底书} \\[1em]
\end{center}

\noindent
\begin{tabular}{|p{3cm}|p{12cm}|}
\hline
公司编号 &  \\
\hline
发明名称 & 基于注意力语义感知的静默KV缓存淘汰与原位压缩的方法、装置、芯片及介质 \\
\hline
申报单位 & 研究院 \\
\hline
申报类型 & 发明 \\
\hline
发明人 & 许达 \\
\hline
技术联系人 & 许达 (xudayj@chinamobile.com, +86-13521894156) \\
\hline
注意事项 & 1．技术联系人应为深入了解本申请提案技术方案的技术人员。 \\
         & 2．请按照集团公司提供的本技术交底书模板逐项填写。 \\
\hline
\end{tabular}

\vspace{1cm}

%==============================================================================
\section{发明名称}
%==============================================================================

基于注意力语义感知的静默KV缓存淘汰与原位压缩的方法、装置、芯片及介质


%==============================================================================
\section{技术领域}
%==============================================================================

本发明涉及人工智能计算、计算机体系结构与存储器管理技术领域，尤其涉及面向包含自注意力机制的神经网络推理（例如 Transformer 大模型推理）的显存/存储层级管理方法与装置，具体为一种在硬件侧旁路获取注意力分数（snoop）并据此实现 KV Cache（键值缓存）静默淘汰（eviction）与原位压缩（in-place compression）的语义感知内存管理单元（Semantic-MMU）。

\textbf{与量子计算的关联说明：}

本发明不涉及量子计算。本字段与其标题沿用模板格式，按模板要求保留。与之对应地，本发明所处的技术背景为：大模型推理正从“算力瓶颈”转向“显存墙与数据调度瓶颈”。在长上下文（10万至百万 token）与多会话并发推理场景下，KV Cache 的容量、带宽、碎片与软件调度开销成为最致命的限制因素。现有 GPU/加速器多数依赖软件实现 KV 的分页管理、压缩与淘汰，导致额外 kernel、同步点与内存搬移，引入显著的尾延迟抖动。本发明采用“算法硅化”思路，将注意力语义信号用于底层存储管理，减少或消除软件参与。

本发明可应用于：
\begin{itemize}
    \item 云端/端侧大模型推理加速器的长上下文推理与多并发推理
    \item 流式对话、RAG 检索增强生成、长文档摘要等长序列推理场景
    \item 需要稳定尾延迟与高吞吐的在线推理服务
    \item 在固定显存/片上 SRAM 约束下提升有效上下文长度的设备
\end{itemize}

\textbf{关联项目：}\underline{\hspace{8cm}}（待填写）

%==============================================================================
\section{现有技术的技术方案}
%==============================================================================

\subsection{量子计算网络安全背景}

本节标题沿用模板。对应到本发明的现有技术背景为：Transformer 推理通常包含计算核心阵列、片上网络（NoC）、片上 SRAM 与外部显存（HBM/GDDR/LPDDR）等层级。推理过程中，KV Cache 为跨 token 的长期状态，其规模随上下文近似线性增长，并在多层多头注意力中被频繁访问。注意力访问呈现显著的动态热点：少数 token 被高频复用，更多历史 token 的注意力贡献长期处于低值。

为控制 KV 的容量与带宽压力，行业常见技术方案包括：
\begin{itemize}
    \item \textbf{分页式/块式 KV 管理}：按页或块分配 KV，软件维护页表与回收策略，降低碎片但增加调度开销；
    \item \textbf{流式与窗口化}：滑动窗口、注意力 sink 等方式丢弃或稀疏化部分历史 token；
    \item \textbf{KV 压缩/量化}：将 KV 从 FP16/BF16 压缩至 INT8/INT4 或采用分组量化减少带宽；
    \item \textbf{软件整理与搬移}：为获得连续空间或提升局部性，触发搬移与重排，带来同步与拷贝。
\end{itemize}

\subsection{后量子密码技术现状}

本节标题沿用模板。对应到本发明的现有技术现状为：上述 KV 管理、淘汰与压缩方案大多由软件在计算核心（GPU SM/NPU 核）与 CPU 之间协同完成。虽然可在一定程度上延长上下文与降低显存压力，但普遍存在下列问题：
\begin{enumerate}
    \item 需要额外 kernel 统计、排序或维护元数据，消耗计算资源并产生同步点；
    \item 淘汰与压缩属于突发性重负载操作，易导致在线推理尾延迟显著波动；
    \item 与推理框架或运行时耦合，跨平台迁移与适配成本高；
    \item 缺乏低成本的硬件语义输入通路，难以在“数据搬运路径”内做出淘汰与压缩决策。
\end{enumerate}

%==============================================================================
\section{现有技术的缺点及本申请提案要解决的技术问题}
%==============================================================================

现有技术存在如下缺陷：

\begin{enumerate}
    \item \textbf{软件侧统计与决策开销大}：注意力分数在计算核心内产生，软件若要基于其进行热度统计与淘汰决策，往往需要额外写回、归并与排序，造成显著开销；
    \item \textbf{对上层不透明或侵入性强}：很多方案需要改动框架/运行时，引入工程风险并限制复用；
    \item \textbf{尾延迟抖动明显}：压缩、碎片整理与回收通常以批处理方式触发，导致抖动；
    \item \textbf{回收与重映射复杂}：软件维护页表/元数据一致性成本高，且常伴随搬移；
    \item \textbf{无法在硬件侧持续衰减}：缺乏硬件级热度衰减计数器与阈值比较，难以实现低成本、持续在线的生命周期管理。
\end{enumerate}

因此，本申请提案要解决的技术问题是：在不改变上层推理框架编程模型的前提下，如何在硬件侧静默地旁路获取注意力语义信息、维护 token 热度衰减计数器，并据此实现 KV Cache 的静默淘汰与原位压缩，同时提供透明重映射与一致性保障，从而显著提升有效上下文长度、吞吐与尾延迟稳定性。

%==============================================================================
\section{本申请提案的技术方案的详细阐述}
%==============================================================================

\subsection{系统架构}

本发明系统包括如下模块：
\begin{enumerate}
    \item \textbf{注意力旁路获取接口（Attention Snoop Tap）}：在片上互连/总线层对注意力分数或其统计量进行旁路捕获，并将其与 token 逻辑标识关联；
    \item \textbf{语义感知内存管理单元（Semantic-MMU）}：部署于内存控制器或片上缓存/显存路径边界，接收语义信号并输出淘汰/压缩/解压控制；
    \item \textbf{Token 热度表与衰减计数器（Heat Table）}：为 token 维护热度计数、衰减、阈值比较与优先级信息；
    \item \textbf{静默淘汰与回收引擎（Silent Evict/Reclaim Engine）}：在后台回收低热度 token 对应的 KV 物理块，维护空闲链表与碎片合并；
    \item \textbf{原位压缩引擎（In-place Compression Engine）}：对中等热度 KV 块执行原位量化/打包存储，并支持按需解压；
    \item \textbf{透明重映射与一致性模块（Transparent Remap \& Consistency）}：维护“语义页表/元数据版本”，对上层提供稳定逻辑视图，确保淘汰、压缩与解压一致性；
    \item \textbf{可观测性计数器与日志缓冲区}：输出淘汰次数、压缩次数、回收字节数、解压命中等指标，支持取证与调参。
\end{enumerate}

\subsection{核心方法步骤}

\subsubsection{1. 分布式密钥生成 (KeyGen)}
本节标题沿用模板。对应到本发明的方法步骤为：在推理会话初始化或上下文窗口建立时，Semantic-MMU 为每个 token 的 KV 块建立元数据记录，包括物理块指针、压缩状态标志、热度计数器初值、版本号、以及可配置阈值（淘汰阈值、压缩阈值、衰减周期与回收预算）。该初始化步骤可由驱动以一次或少量配置写入完成，后续运行无需软件持续介入。

\subsubsection{2. 消息哈希与映射}
本节标题沿用模板。对应到本发明的方法步骤为：Semantic-MMU 对注意力旁路信号进行标识映射。具体地，将注意力分数（或聚合统计量，例如跨头最大值、top-k 贡献和、滑动均值等）与 token 的逻辑标识（layer/head/token index/session id）绑定，并映射到 Heat Table 表项。映射可采用直接索引或哈希分桶，以降低元数据存储规模与冲突。

\subsubsection{3. 分布式高斯采样}
本节标题沿用模板。对应到本发明的方法步骤为：Semantic-MMU 对热度计数器执行持续衰减与在线更新。可采用指数衰减与定点近似实现，例如：
\[
H \leftarrow \alpha \cdot H + \beta \cdot S
\]
其中 $S$ 为旁路语义贡献，$\alpha \in (0,1)$ 为衰减系数，$\beta$ 为增益系数。为避免热度长期“锁死”，还可引入概率式采样更新，使冷 token 更快进入淘汰/压缩候选集合，且硬件开销保持极低。

\subsubsection{4. 基于 NTT 的线性变换}
本节标题沿用模板。对应到本发明的方法步骤为：当热度低于压缩阈值但高于淘汰阈值时，触发原位压缩引擎对 KV 块执行表示转换与打包存储，例如分组量化、缩放因子生成、离群值处理与 packing。该过程抽象为：
\[
(\mathbf{K},\mathbf{V}) \rightarrow \text{Pack}(\text{Quant}(\mathbf{K}), \text{Quant}(\mathbf{V}), \text{Scale})
\]
并在同一物理块或预留邻接块中就地写回，避免生成临时张量与大规模搬移。

\subsubsection{5. 基于 Beaver 三元组的安全范数验证}
本节标题沿用模板。对应到本发明的方法步骤为：为保证压缩与淘汰不会破坏推理正确性与一致性，Semantic-MMU 维护轻量校验元数据与误差界信息，例如 CRC/校验和、量化误差上界估计、缩放因子范围、版本号等，并在解压或关键访问路径执行快速一致性检查。若校验失败或误差界超出策略阈值，则对该 KV 块执行回退策略（例如强制保留高精度、禁止再次压缩、提升热度），以保证系统稳定性。

\subsubsection{6. 动态节点管理}
本节标题沿用模板。对应到本发明的方法步骤为：在运行期，当显存压力达到阈值或回收预算触发时，Semantic-MMU 在后台静默选择低热度 token 的 KV 物理块执行淘汰：释放物理块、维护空闲链表、合并空洞并更新语义页表；对仍需保留的历史 token，可保留压缩态或下沉到更低速存储层。整个过程不触发软件中断，不引入显式同步点，对上层框架保持透明。

%==============================================================================
\section{本申请提案的关键点和欲保护点}
%==============================================================================

本申请提案主要包含以下关键创新点和保护点：

\begin{enumerate}
    \item \textbf{注意力语义旁路获取（Bus/NoC Snoop）}：在不增加计算核心负担的情况下旁路捕获注意力分数/统计量，并与 token 元数据绑定；
    \item \textbf{硬件级热度衰减计数器}：硬件侧持续维护 token 热度并进行阈值比较，输出淘汰与压缩决策；
    \item \textbf{静默淘汰与回收状态机}：在后台完成物理块回收与碎片合并，不阻塞主流水线；
    \item \textbf{原位压缩与按需解压}：对 KV 块执行原位量化打包，并在访问路径按需解压，降低显存与带宽；
    \item \textbf{透明重映射与一致性保障}：维护语义页表/版本控制，使上层以为 KV 逻辑地址稳定存在；
    \item \textbf{可配置策略与可观测性}：阈值、预算、优先级与回退策略可配置，并输出计数器用于调参与取证。
\end{enumerate}

%==============================================================================
\section{与第三条中最接近的现有技术相比，本申请提案有何技术优点}
%==============================================================================

相比于现有技术（如软件侧分页式 KV 管理、软件淘汰与软件压缩/整理等），本方案具有显著优势：

(1) \textbf{对上层框架透明}：语义页表与元数据在硬件侧维护，上层无需引入额外管理逻辑。
(2) \textbf{显著降低软件开销与尾延迟抖动}：淘汰回收与压缩在硬件后台执行，避免突发性软件整理。
(3) \textbf{提高显存利用率与有效上下文长度}：低热度 KV 被回收，中等热度 KV 原位压缩，固定显存下可显著扩大可支持上下文。
(4) \textbf{带宽效率更高}：压缩态降低 KV 读写带宽，热度驱动减少无效历史 token 访问。
(5) \textbf{易落地}：不依赖新材料与新工艺，主要为内存控制器、互连旁路与少量定点计算/状态机逻辑，可用成熟 CMOS 数字逻辑实现。

%==============================================================================
\section{其他有助于理解本申请提案的技术资料}
%==============================================================================

(1) vLLM: PagedAttention（公开论文/开源实现）
(2) FlashAttention / FlashAttention-2（公开论文/开源实现）
(3) StreamingLLM / Attention Sinks（公开论文/开源实现）
(4) 统一虚拟内存（UVM）与内存压缩相关公开资料（GPU/加速器白皮书与学术论文）

%==============================================================================
\section{本申请提案的侵权证据可获得性}
%==============================================================================

本提案的侵权行为具有较高的可检测性：
(1) \textbf{硬件接口与寄存器/计数器可观测性}：实现 Semantic-MMU 的加速器通常具有阈值、预算、策略等可配置寄存器，以及淘汰/压缩/解压等性能计数器，可通过驱动/SDK 识别。
(2) \textbf{微架构结构性证据}：芯片/固件分析中若存在注意力分数旁路通路、token 热度表、后台回收状态机与原位量化打包逻辑，可形成强结构证据。
(3) \textbf{行为侧证（性能曲线）}：在长上下文压力测试下呈现“显存接近满载但尾延迟稳定、且 KV 压缩态切换与带宽曲线相关”的特征，可作为侵权线索。


\end{document}
